\section*{Glosario de Acrónimos}
\markboth{Glosario}{Glosario}

\begin{description}

  \item[API] \textit{Application Programming Interface}. Interfaz de programación de aplicaciones; conjunto de reglas y protocolos que permiten la comunicación entre sistemas de software.

  \item[BI] \textit{Business Intelligence}. Conjunto de estrategias, procesos y herramientas orientadas a la transformación de datos en información de apoyo a la toma de decisiones empresariales.

  \item[CAIQ] \textit{Career Alignment and Insight Qualifier}. Sistema de orientación profesional desarrollado en este trabajo que analiza el perfil técnico de un candidato, detecta su \textit{skill gap} respecto a un rol objetivo y genera recomendaciones personalizadas de formación y empleo.

  \item[CSV] \textit{Comma-Separated Values}. Formato de archivo de texto plano para almacenar datos tabulares separados por comas.

  \item[CPU] \textit{Central Processing Unit}. Unidad central de procesamiento; componente principal de un ordenador responsable de ejecutar instrucciones.

  \item[CV] \textit{Curriculum Vitae}. Documento que recoge la trayectoria académica y profesional de una persona. En este trabajo constituye la entrada principal del sistema de recomendación.

  \item[DOCX] Formato de archivo propietario de Microsoft Word basado en el estándar Office Open XML.

  \item[EDA] \textit{Exploratory Data Analysis} (Análisis Exploratorio de Datos). Fase del proceso analítico orientada a caracterizar la estructura, calidad y distribución de los datos antes de aplicar modelos.

  \item[ETL] \textit{Extract, Transform, Load}. Proceso de extracción, transformación y carga de datos desde fuentes heterogéneas hacia un repositorio estructurado y homogéneo.

  \item[F1] \textit{F1-score}. Media armónica de la precisión y el \textit{recall}; métrica de evaluación de clasificadores binarios.

  \item[FN] \textit{False Negative}. Instancia positiva clasificada erróneamente como negativa.

  \item[FP] \textit{False Positive}. Instancia negativa clasificada erróneamente como positiva.

  \item[GPU] \textit{Graphics Processing Unit}. Unidad de procesamiento gráfico; utilizada en aprendizaje profundo por su capacidad de cómputo paralelo masivo.

  \item[JSON] \textit{JavaScript Object Notation}. Formato ligero de intercambio de datos basado en texto, ampliamente utilizado en APIs y archivos de configuración.

  \item[LLM] \textit{Large Language Model}. Modelo de lenguaje de gran escala entrenado sobre grandes volúmenes de texto; base de sistemas como GPT o LLaMA.

  \item[ML] \textit{Machine Learning} (Aprendizaje Automático). Rama de la inteligencia artificial que desarrolla algoritmos capaces de aprender patrones a partir de datos.

  \item[MLOps] \textit{Machine Learning Operations}. Conjunto de prácticas y herramientas para el despliegue, monitorización y mantenimiento de modelos de aprendizaje automático en producción.

  \item[NDCG] \textit{Normalized Discounted Cumulative Gain}. Métrica de evaluación de sistemas de recuperación de información y recomendación que pondera la relevancia de los resultados por su posición en el ranking.

  \item[NER] \textit{Named Entity Recognition} (Reconocimiento de Entidades Nombradas). Tarea de procesamiento del lenguaje natural orientada a identificar y clasificar entidades (personas, organizaciones, fechas, etc.) en texto.

  \item[NLP] \textit{Natural Language Processing} (Procesamiento del Lenguaje Natural). Rama de la inteligencia artificial que estudia la interacción entre ordenadores y lenguaje humano.

  \item[PDF] \textit{Portable Document Format}. Formato de archivo desarrollado por Adobe Systems para representar documentos de forma independiente del \textit{software} y el \textit{hardware}.

  \item[SQL] \textit{Structured Query Language}. Lenguaje estándar para la gestión y consulta de bases de datos relacionales.

  \item[TFM] Trabajo Fin de Máster. Proyecto académico de investigación o desarrollo que integra las competencias adquiridas durante un máster universitario.

  \item[TP] \textit{True Positive}. Instancia positiva clasificada correctamente como positiva.

  \item[TXT] Formato de archivo de texto plano sin formato adicional.

  \item[URL] \textit{Uniform Resource Locator}. Dirección que identifica de forma unívoca un recurso en Internet.

  \item[XML] \textit{Extensible Markup Language}. Lenguaje de marcado diseñado para almacenar y transportar datos de forma estructurada y legible por humanos y máquinas.

\end{description}
